% sec2_4.tex — Section 2.4: Hypothesis Testing

Statistical inference in large-sample theory is based on test statistics whose
asymptotic distributions are known under the truth of the null hypothesis.
Derivation of the distribution of test statistics is easier than in finite-sample
theory because we are only concerned about the large-sample \emph{approximation}
to the exact distribution.  In this section we derive test statistics, assuming
throughout that a consistent estimator, $\widehat{\mathbf{S}}$, of
$\mathbf{S}$ $(\equiv \E(\mathbf{g}_i\mathbf{g}_i'))$ is available.  The issue of
consistent estimation of $\mathbf{S}$ will be taken up in the next section.

%% ─── Testing Linear Hypotheses ─────────────────────────────────────────────
\subsection*{Testing Linear Hypotheses}

Consider testing a hypothesis about the $k$-th coefficient $\beta_k$.
Proposition~\ref{prop:2.1} implies that under the null
$\mathrm{H}_0\colon \beta_k = \bar{\beta}_k$,
\[
  \sqrt{n}(b_k - \bar{\beta}_k) \dto N\!\big(0,\, \avar(b_k)\big)
  \quad \text{and} \quad
  \widehat{\avar(b_k)} \pto \avar(b_k),
\]
where $b_k$ is the $k$-th element of $\mathbf{b}$ and $\avar(b_k)$ is the $(k,k)$
element of the $K \times K$ matrix $\avar(\mathbf{b})$.  So Lemma~\ref{lem:2.4}(c)
guarantees that
\begin{equation}
  t_k \equiv \frac{\sqrt{n}(b_k - \bar{\beta}_k)}
  {\sqrt{\widehat{\avar(b_k)}}}
  = \frac{b_k - \bar{\beta}_k}{SE^*(b_k)}
  \dto N(0, 1),
  \label{eq:2.4.1}
\end{equation}
where
\[
  SE^*(b_k) \equiv \sqrt{\tfrac{1}{n} \cdot \widehat{\avar(b_k)}}
  \equiv \sqrt{\tfrac{1}{n} \cdot
  (\mathbf{S}_{\mathbf{xx}}^{-1}\,\widehat{\mathbf{S}}\,
  \mathbf{S}_{\mathbf{xx}}^{-1})_{kk}}.
\]
The denominator in this $t$-ratio, $SE^*(b_k)$, is called the
\textbf{heteroskedasticity-consistent standard error},
\textbf{(heteroskedasticity-)robust standard error}, or \textbf{White's standard
error}.  The reason for this terminology is that the error term can be conditionally
heteroskedastic; recall that we have not assumed conditional homoskedasticity (that
$\E(\varepsilon_i^2 \mid \mathbf{x}_i)$ does not depend on $\mathbf{x}_i$) to derive
the asymptotic distribution of $t_k$.  This $t$-ratio is called the \textbf{robust}
$t$-ratio, to distinguish it from the $t$-ratio of Chapter~1.  The relationship
between these two sorts of $t$-ratio will be discussed in Section~2.6.

Given this $t$-ratio, testing the null hypothesis
$\mathrm{H}_0\colon \beta_k = \bar{\beta}_k$ at a significance level of $\alpha$
proceeds as follows:

\smallskip
\noindent\emph{Step~1:} Calculate $t_k$ by the formula \eqref{eq:2.4.1}.

\noindent\emph{Step~2:} Look up the table of $N(0,1)$ to find the critical value
$t_{\alpha/2}$ which leaves $\alpha/2$ to the upper tail of the standard normal
distribution.  (example: if $\alpha = 5\%$, $t_{\alpha/2} = 1.96$.)

\noindent\emph{Step~3:} Accept the hypothesis if $|t_k| < t_{\alpha/2}$; otherwise
reject.

\smallskip
The differences from the finite-sample $t$-test are: (1) the way the standard error
is calculated is different, (2) we use the table of $N(0,1)$ rather than that of
$t(n-K)$, and (3) the \textbf{actual size} or \textbf{exact size} of the test (the
probability of Type~I error given the sample size) equals the \textbf{nominal size}
(i.e., the desired significance level $\alpha$) only approximately, although the
approximation becomes arbitrarily good as the sample size increases.  The difference
between the exact size and the nominal size of a test is called the \textbf{size
distortion}.  Since $t_k$ is asymptotically standard normal, the size distortion of
the $t$-test converges to zero as the sample size $n$ goes to infinity.

Thus, we have proved the first half of

\begin{proposition}[robust $t$-ratio and Wald statistic]
\label{prop:2.3}
Suppose Assumptions~\ref{ass:2.1}--\ref{ass:2.5} hold, and suppose there is
available a consistent estimate $\widehat{\mathbf{S}}$ of\/
$\mathbf{S}$ $(= \E(\mathbf{g}_i\mathbf{g}_i'))$.  As before, let
\[
  \widehat{\avar(\mathbf{b})}
  \equiv \mathbf{S}_{\mathbf{xx}}^{-1}\,\widehat{\mathbf{S}}\,
  \mathbf{S}_{\mathbf{xx}}^{-1}.
\]
Then
\begin{enumerate}[label=(\alph*)]
  \item Under the null hypothesis $\mathrm{H}_0\colon \beta_k = \bar{\beta}_k$,
    $t_k$ defined in \eqref{eq:2.4.1} $\dto N(0,1)$.
  \item Under the null hypothesis
    $\mathrm{H}_0\colon \mathbf{R}\bm{\beta} = \mathbf{r}$, where $\mathbf{R}$ is an
    $\#\mathbf{r} \times K$ matrix (where $\#\mathbf{r}$, the dimension of\/
    $\mathbf{r}$, is the number of restrictions) of full row rank,
    \begin{equation}
      W \equiv n \cdot (\mathbf{R}\mathbf{b} - \mathbf{r})'
      \{\mathbf{R}[\widehat{\avar(\mathbf{b})}]\mathbf{R}'\}^{-1}
      (\mathbf{R}\mathbf{b} - \mathbf{r})
      \dto \chi^2(\#\mathbf{r}).
      \label{eq:2.4.2}
    \end{equation}
\end{enumerate}
\end{proposition}

What remains to be shown is that $W \dto \chi^2(\#\mathbf{r})$, which is a
straightforward application of Lemma~\ref{lem:2.4}(d).

\begin{proof}[Proof (continued)]
Write $W$ as
\[
  W = \mathbf{c}_n'\mathbf{Q}_n^{-1}\mathbf{c}_n
  \quad \text{where} \quad
  \mathbf{c}_n \equiv \sqrt{n}(\mathbf{R}\mathbf{b} - \mathbf{r})
  \;\text{and}\;
  \mathbf{Q}_n \equiv \mathbf{R}\,\widehat{\avar(\mathbf{b})}\,\mathbf{R}'.
\]
Under $\mathrm{H}_0$,
$\mathbf{c}_n = \mathbf{R}\sqrt{n}(\mathbf{b} - \bm{\beta})$.  So by
Proposition~\ref{prop:2.1},
\[
  \mathbf{c}_n \dto \mathbf{c} \quad \text{where} \quad
  \mathbf{c} \sim N\!\big(\mathbf{0},\,
  \mathbf{R}\,\avar(\mathbf{b})\,\mathbf{R}'\big).
\]
Also by Proposition~\ref{prop:2.1},
\[
  \mathbf{Q}_n \pto \mathbf{Q} \quad \text{where} \quad
  \mathbf{Q} \equiv \mathbf{R}\,\avar(\mathbf{b})\,\mathbf{R}'.
\]
Because $\mathbf{R}$ is of full row rank and $\avar(\mathbf{b})$ is positive
definite, $\mathbf{Q}$ is invertible.  Therefore, by Lemma~\ref{lem:2.4}(d),
\[
  W \dto \mathbf{c}'\mathbf{Q}^{-1}\mathbf{c}.
\]
Since the $\#\mathbf{r}$-dimensional random vector $\mathbf{c}$ is normally
distributed and since $\mathbf{Q}$ equals $\Var(\mathbf{c})$,
$\mathbf{c}'\mathbf{Q}^{-1}\mathbf{c} \sim \chi^2(\#\mathbf{r})$.
\end{proof}

This chi-square statistic $W$ is a Wald statistic because it is based on unrestricted
estimates ($\mathbf{b}$ and $\widehat{\avar(\mathbf{b})}$ here) not constrained by
the null hypothesis $\mathrm{H}_0$.  Testing $\mathrm{H}_0$ at a significance level
of $\alpha$ proceeds as follows.

\smallskip
\noindent\emph{Step~1:} Calculate the $W$ statistic by the formula \eqref{eq:2.4.2}.

\noindent\emph{Step~2:} Look up the table of $\chi^2(\#\mathbf{r})$ distribution to
find the critical value $\chi_\alpha^2(\#\mathbf{r})$ that gives $\alpha$ to the upper
tail of the $\chi^2(\#\mathbf{r})$ distribution.

\noindent\emph{Step~3:} If $W < \chi_\alpha^2(\#\mathbf{r})$, then accept
$\mathrm{H}_0$; otherwise reject.

\smallskip
The probability of Type~I error approaches $\alpha$ as the sample becomes larger.
As will be made clear in Section~2.6, this Wald statistic is closely related to the
familiar $F$-test under conditional homoskedasticity.

%% ─── The Test Is Consistent ───────────────────────────────────────────────
\subsection*{The Test Is Consistent}

Recall from basic statistics that the (finite-sample) \textbf{power} of a test is the
probability of rejecting the null hypothesis when it is false given a sample of finite
size (that is, the power is 1 minus the probability of Type~II error).  Power will
obviously depend on the DGP (i.e., how the data were actually generated) considered
as the alternative as well as on the size (significance level) of the test.  For
example, consider any DGP $\{y_i, \mathbf{x}_i\}$ satisfying
Assumptions~\ref{ass:2.1}--\ref{ass:2.5} but \emph{not} the null hypothesis
$\mathrm{H}_0\colon \beta_k = \bar{\beta}_k$.  The power of the $t$-test of size
$\alpha$ against this alternative is
\[
  \text{power} = \mathrm{Prob}(|t_k| > t_{\alpha/2}),
\]
which depends on the DGP in question because the DGP controls the distribution of
$t_k$.  We say that a test is \textbf{consistent} against a set of DGPs, none of which
satisfies the null, if the power against any particular member of the set approaches
unity as $n \to \infty$ for any assumed significance level.

That the $t$-test is consistent against the set of alternatives (DGPs) satisfying
Assumptions~\ref{ass:2.1}--\ref{ass:2.5} can be seen as follows.  Look at the
expression \eqref{eq:2.4.1} for the $t$-ratio, reproduced here:
\[
  t_k \equiv \frac{\sqrt{n}(b_k - \bar{\beta}_k)}
  {\sqrt{\widehat{\avar(b_k)}}}.
\]
The denominator converges to $\sqrt{\avar(b_k)}$ despite the fact that the DGP does
not satisfy the null (recall that all parts of Proposition~\ref{prop:2.1} hold
regardless of the truth of the null, provided Assumptions~\ref{ass:2.1}--\ref{ass:2.5}
are satisfied).  On the other hand, the numerator tends to $+\infty$ or $-\infty$
because $b_k$ converges in probability to the DGP's $\beta_k$, which is different
from $\bar{\beta}_k$.  So the power tends to unity as the sample size $n$ tends to
infinity, implying that the $t$-test of Proposition~\ref{prop:2.3} is consistent
against those alternatives, the DGPs that do not satisfy the null.  The same is true
for the Wald test.

%% ─── Asymptotic Power ─────────────────────────────────────────────────────
\subsection*{Asymptotic Power}

For later use in the next chapter, we define here the \textbf{asymptotic power} of a
consistent test.  As noted above, the power of the $t$-test approaches to unity as
the sample size increases while the DGP taken as the alternative is held fixed.
But if the DGP gets closer and closer to the null as the sample size increases, the
power may not converge to unity.  A sequence of such DGPs is called a
\textbf{sequence of local alternatives}.  For the regression model and for the null
of $\mathrm{H}_0\colon \beta_k = \bar{\beta}_k$, it is a sequence of DGPs such
that (i) the $n$-th DGP, $\{y_i^{(n)}, \mathbf{x}_i^{(n)}\}$
$(i = 1, 2, \ldots)$, satisfies Assumptions~\ref{ass:2.1}--\ref{ass:2.5} and
converges in a certain sense to a fixed DGP
$\{y_i, \mathbf{x}_i\}$,\footnote{See, e.g., Assumption~1 of Newey (1985) for a
precise statement.} and (ii)~the value of $\beta_k$ of the $n$-th DGP,
$\beta_k^{(n)}$, converges to $\bar{\beta}_k$.  Suppose, further, that
$\beta_k^{(n)}$ satisfies
\begin{equation}
  \beta_k^{(n)} = \bar{\beta}_k + \frac{\gamma}{\sqrt{n}}
  \label{eq:2.4.3}
\end{equation}
for some given $\gamma \neq 0$.  So $\beta_k^{(n)}$ approaches to $\bar{\beta}_k$
at a rate proportional to $1/\sqrt{n}$.  This special sequence of local alternatives
is called a \textbf{Pitman drift} or a \textbf{Pitman sequence}.  Substituting
\eqref{eq:2.4.3} into \eqref{eq:2.4.1}, the $t$-ratio above can be rewritten as
\begin{equation}
  t_k = \frac{\sqrt{n}(b_k - \beta_k^{(n)})}
  {\sqrt{\widehat{\avar(b_k)}}}
  + \frac{\gamma}{\sqrt{\widehat{\avar(b_k)}}}.
  \label{eq:2.4.4}
\end{equation}

If the sample of size $n$ is generated by the $n$-th DGP of a Pitman drift, does
$t_k$ converge to a nontrivial distribution?  Since the $n$-th DGP satisfies
Assumptions~\ref{ass:2.1}--\ref{ass:2.5}, the first term on the right hand side of
\eqref{eq:2.4.4} converges in distribution to $N(0,1)$ by parts~(b) and (c) of
Proposition~\ref{prop:2.1}.  By part~(c) of Proposition~\ref{prop:2.1} and the fact
that $\{y_i^{(n)}, \mathbf{x}_i^{(n)}\}$ ``converges'' to a fixed DGP, the second
term converges in probability to
\begin{equation}
  \mu \equiv \frac{\gamma}{\sqrt{\avar(b_k)}},
  \label{eq:2.4.5}
\end{equation}
where $\avar(b_k)$ is evaluated at the fixed DGP.  Therefore,
$t_k \dto N(\mu, 1)$ along this sequence of local alternatives.  If the significance
level is $\alpha$, the power converges to
\begin{equation}
  \mathrm{Prob}(|x| > t_{\alpha/2})
  \label{eq:2.4.6}
\end{equation}
where $x \sim N(\mu, 1)$ and $t_{\alpha/2}$ is the level-$\alpha$ critical value.
This probability is called the \textbf{asymptotic power}.  It is a measure of the
ability of the test to detect small deviations of the model from the null hypothesis.
Evidently, the larger is $|\mu|$, the higher is the asymptotic power for any given
size $\alpha$.  By a similar argument, it is easy to show that the Wald statistic
converges to a distribution called \textbf{noncentral chi-squared}.

%% ─── Testing Nonlinear Hypotheses ─────────────────────────────────────────
\subsection*{Testing Nonlinear Hypotheses}

The Wald statistic can be generalized to a test of a set of nonlinear restrictions on
$\bm{\beta}$.  Consider a null hypothesis of the form
\[
  \mathrm{H}_0\colon \mathbf{a}(\bm{\beta}) = \mathbf{0}.
\]
Here, $\mathbf{a}$ is a vector-valued function with continuous first derivatives.
Let $\#\mathbf{a}$ be the dimension of $\mathbf{a}(\bm{\beta})$ (so the null
hypothesis has $\#\mathbf{a}$ restrictions), and let $\mathbf{A}(\bm{\beta})$ be the
$\#\mathbf{a} \times K$ matrix of first derivatives evaluated at $\bm{\beta}$:
$\mathbf{A}(\bm{\beta}) = \partial \mathbf{a}(\bm{\beta})/\partial \bm{\beta}'$.
For the hypothesis to be well-defined, we assume that $\mathbf{A}(\bm{\beta})$ is
of full row rank (this is the generalization of the requirement for linear hypothesis
$\mathbf{R}\bm{\beta} = \mathbf{r}$ that $\mathbf{R}$ is of full row rank).
Lemma~\ref{lem:2.5} of Section~2.1 and Proposition~\ref{prop:2.1}(b) imply that
\begin{equation}
  \sqrt{n}[\mathbf{a}(\mathbf{b}) - \mathbf{a}(\bm{\beta})]
  \dto \mathbf{c}, \quad
  \mathbf{c} \sim N\!\big(\mathbf{0},\,
  \mathbf{A}(\bm{\beta})\,\avar(\mathbf{b})\,\mathbf{A}(\bm{\beta})'\big).
  \label{eq:2.4.7}
\end{equation}
Since $\mathbf{a}(\bm{\beta}) = \mathbf{0}$ under $\mathrm{H}_0$,
\eqref{eq:2.4.7} becomes
\begin{equation}
  \sqrt{n}\,\mathbf{a}(\mathbf{b}) \dto \mathbf{c}, \quad
  \mathbf{c} \sim N\!\big(\mathbf{0},\,
  \mathbf{A}(\bm{\beta})\,\avar(\mathbf{b})\,\mathbf{A}(\bm{\beta})'\big).
  \label{eq:2.4.8}
\end{equation}
Since $\mathbf{b} \pto \bm{\beta}$ by Proposition~\ref{prop:2.1}(a),
Lemma~\ref{lem:2.3}(a) implies that
$\mathbf{A}(\mathbf{b}) \pto \mathbf{A}(\bm{\beta})$.  By
Proposition~\ref{prop:2.1}(c),
$\widehat{\avar(\mathbf{b})} \pto \avar(\mathbf{b})$.  So by
Lemma~\ref{lem:2.3}(a),
\begin{equation}
  \mathbf{A}(\mathbf{b})\,\widehat{\avar(\mathbf{b})}\,\mathbf{A}(\mathbf{b})'
  \pto \mathbf{A}(\bm{\beta})\,\avar(\mathbf{b})\,\mathbf{A}(\bm{\beta})'
  = \Var(\mathbf{c}).
  \label{eq:2.4.9}
\end{equation}
Because $\mathbf{A}(\bm{\beta})$ is of full row rank and $\avar(\mathbf{b})$ is
positive definite, $\Var(\mathbf{c})$ is invertible.  Then Lemma~\ref{lem:2.4}(d),
\eqref{eq:2.4.8}, and \eqref{eq:2.4.9} imply
\begin{equation}
  \sqrt{n}\,\mathbf{a}(\mathbf{b})'
  \{\mathbf{A}(\mathbf{b})\,\widehat{\avar(\mathbf{b})}\,
  \mathbf{A}(\mathbf{b})'\}^{-1}
  \sqrt{n}\,\mathbf{a}(\mathbf{b})
  \dto \mathbf{c}'\,\Var(\mathbf{c})^{-1}\,\mathbf{c}
  \sim \chi^2(\#\mathbf{a}).
  \label{eq:2.4.10}
\end{equation}
Combining two $\sqrt{n}$'s in \eqref{eq:2.4.10} into one $n$, we have proved

\paragraph{Proposition 2.3 (continued):}

\begin{enumerate}[label=(\alph*), start=3]
  \item \emph{Under the null hypothesis with $\#\mathbf{a}$ restrictions
    $\mathrm{H}_0\colon \mathbf{a}(\bm{\beta}) = \mathbf{0}$, such that
    $\mathbf{A}(\bm{\beta})$, the $\#\mathbf{a} \times K$ matrix of continuous first
    derivatives of\/ $\mathbf{a}(\bm{\beta})$, is of full row rank, we have}
    \begin{equation}
      W \equiv n \cdot \mathbf{a}(\mathbf{b})'
      \{\mathbf{A}(\mathbf{b})\,\widehat{\avar(\mathbf{b})}\,
      \mathbf{A}(\mathbf{b})'\}^{-1}
      \mathbf{a}(\mathbf{b})
      \dto \chi^2(\#\mathbf{a}).
      \label{eq:2.4.11}
    \end{equation}
\end{enumerate}

Part~(c) is a generalization of (b); by setting
$\mathbf{a}(\bm{\beta}) = \mathbf{R}\bm{\beta} - \mathbf{r}$,
\eqref{eq:2.4.11} reduces to \eqref{eq:2.4.2}, the Wald statistic for linear
restrictions.

The choice of $\mathbf{a}(\cdot)$ for representing a given set of restrictions is not
unique.  For example, $\beta_1\beta_2 = 1$ can be written as
$a(\bm{\beta}) = 0$ with $a(\bm{\beta}) = \beta_1\beta_2 - 1$ or with
$a(\bm{\beta}) = \beta_1 - 1/\beta_2$.  While part~(c) of the proposition guarantees
that in large samples the outcome of the Wald test is the same regardless of the
choice of the function $\mathbf{a}$, the numerical value of the Wald statistic $W$
does depend on the representation, and the test outcome can be different in finite
samples.  In the above example, the second representation,
$a(\bm{\beta}) = \beta_1 - 1/\beta_2$, does not satisfy the requirement of continuous
derivatives at $\beta_2 = 0$.  Indeed, a Monte Carlo study by Gregory and Veall
(1985) reports that, when $\beta_2$ is close to zero, the Wald test based on the
second representation rejects the null too often in small samples.
